\documentclass[10pt,a4paper]{amsart}
\usepackage{amsmath,amssymb,amsthm,graphicx}
% Retained preamble for the Stacks Project chapter morphisms.tex (Morphisms of Schemes).
% The chapter file pinned-original-source.tex carries only \input{preamble} on its first line, so no
% preamble text lives inside the chapter itself.  The Stacks Project uses one single shared preamble for
% every chapter; the snapshot stages it verbatim in sources/stacks-chow/preamble.tex.  The block below is
% copied mechanically, line for line, from that file's lines 14-18 (the Xy-pic setup) and lines 154-217
% (the native theorem declarations together with the mathematical macros).  Nothing was added, removed or
% redefined.  The wrapper class and the amsmath/amssymb/amsthm/graphicx base packages are supplied by the
% assembler, so the class/\externaldocument part of the shared preamble is not repeated here.
% For commutative diagrams we use Xy-pic
\usepackage[all]{xy}

% We use 2cell for 2-commutative diagrams.
\xyoption{2cell}
% Theorem environments.
%
\theoremstyle{plain}
\newtheorem{theorem}[subsection]{Theorem}
\newtheorem{proposition}[subsection]{Proposition}
\newtheorem{lemma}[subsection]{Lemma}

\theoremstyle{definition}
\newtheorem{definition}[subsection]{Definition}
\newtheorem{example}[subsection]{Example}
\newtheorem{exercise}[subsection]{Exercise}
\newtheorem{situation}[subsection]{Situation}

\theoremstyle{remark}
\newtheorem{remark}[subsection]{Remark}
\newtheorem{remarks}[subsection]{Remarks}

\numberwithin{equation}{subsection}

% Macros
%
\def\lim{\mathop{\mathrm{lim}}\nolimits}
\def\colim{\mathop{\mathrm{colim}}\nolimits}
\def\Spec{\mathop{\mathrm{Spec}}}
\def\Hom{\mathop{\mathrm{Hom}}\nolimits}
\def\Ext{\mathop{\mathrm{Ext}}\nolimits}
\def\SheafHom{\mathop{\mathcal{H}\!\mathit{om}}\nolimits}
\def\SheafExt{\mathop{\mathcal{E}\!\mathit{xt}}\nolimits}
\def\Sch{\mathit{Sch}}
\def\Mor{\mathop{\mathrm{Mor}}\nolimits}
\def\Ob{\mathop{\mathrm{Ob}}\nolimits}
\def\Sh{\mathop{\mathit{Sh}}\nolimits}
\def\NL{\mathop{N\!L}\nolimits}
\def\CH{\mathop{\mathrm{CH}}\nolimits}
\def\proetale{{pro\text{-}\acute{e}tale}}
\def\etale{{\acute{e}tale}}
\def\QCoh{\mathit{QCoh}}
\def\Ker{\mathop{\mathrm{Ker}}}
\def\Im{\mathop{\mathrm{Im}}}
\def\Coker{\mathop{\mathrm{Coker}}}
\def\Coim{\mathop{\mathrm{Coim}}}

% Boxtimes
%
\DeclareMathSymbol{\boxtimes}{\mathbin}{AMSa}{"02}

%
% Macros for moduli stacks/spaces
%
\def\QCohstack{\mathcal{QC}\!\mathit{oh}}
\def\Cohstack{\mathcal{C}\!\mathit{oh}}
\def\Spacesstack{\mathcal{S}\!\mathit{paces}}
\def\Quotfunctor{\mathrm{Quot}}
\def\Hilbfunctor{\mathrm{Hilb}}
\def\Curvesstack{\mathcal{C}\!\mathit{urves}}
\def\Polarizedstack{\mathcal{P}\!\mathit{olarized}}
\def\Complexesstack{\mathcal{C}\!\mathit{omplexes}}
% \Pic is the operator that assigns to X its picard group, usage \Pic(X)
% \Picardstack_{X/B} denotes the Picard stack of X over B
% \Picardfunctor_{X/B} denotes the Picard functor of X over B
\def\Pic{\mathop{\mathrm{Pic}}\nolimits}
\def\Picardstack{\mathcal{P}\!\mathit{ic}}
\def\Picardfunctor{\mathrm{Pic}}
\def\Deformationcategory{\mathcal{D}\!\mathit{ef}}

% --- appended 2026-10-12 (Chief): the project own definitions of the reference/slogan/history
% display environments, copied mechanically from the shared Stacks preamble lines 8-12
% (\usepackage{verbatim} plus the three \newenvironment lines, which the chapter preamble keeps
% verbatim from the project preamble). Nothing added, removed or redefined: these environments
% carry no mathematical content and the project itself typesets them as comments. Without these
% three lines a counted statement whose own text uses them fails with Environment slogan undefined.
\usepackage{verbatim}
\newenvironment{reference}{\comment}{\endcomment}
\newenvironment{slogan}{\comment}{\endcomment}
\newenvironment{history}{\comment}{\endcomment}
\usepackage{cleveref}
\crefname{originallemma}{original lemma}{original lemmas}
\Crefname{originallemma}{Original lemma}{Original lemmas}
\makeatletter
\newlabel{lemma-characterize-proper}{{0BX5}{source}{}{sourcepointer-1}{}}
\newlabel{lemma-characterize-proper@cref}{{[originallemma][0][]0BX5}{[1][0][]source}{}{}{}}
\makeatother
\crefname{originallemma}{original lemma}{original lemmas}
\Crefname{originallemma}{Original lemma}{Original lemmas}
\makeatletter
\newlabel{lemma-morphism-defined-local-ring}{{0BX6}{source}{}{sourcepointer-2}{}}
\newlabel{lemma-morphism-defined-local-ring@cref}{{[originallemma][0][]0BX6}{[1][0][]source}{}{}{}}
\makeatother
\crefname{originallemma}{original lemma}{original lemmas}
\Crefname{originallemma}{Original lemma}{Original lemmas}
\makeatletter
\newlabel{properties-lemma-ring-affine-open-injective-local-ring}{{0BX3}{source}{}{sourcepointer-3}{}}
\newlabel{properties-lemma-ring-affine-open-injective-local-ring@cref}{{[originallemma][0][]0BX3}{[1][0][]source}{}{}{}}
\makeatother
\crefname{originallemma}{original lemma}{original lemmas}
\Crefname{originallemma}{Original lemma}{Original lemmas}
\makeatletter
\newlabel{schemes-lemma-where-are-they-equal}{{01KM}{source}{}{sourcepointer-4}{}}
\newlabel{schemes-lemma-where-are-they-equal@cref}{{[originallemma][0][]01KM}{[1][0][]source}{}{}{}}
\makeatother
\title{A morphism from an open subscheme into a proper scheme over the base extends to an open neighbourhood of a valuation-theoretic point in the closure}
\author{The Stacks Project Authors}
\date{stacks-project:morphisms.tex:local-snapshot-20261012; GNU FDL 1.2 or later}
\begin{document}
\maketitle
\noindent\textit{Editorial note.} All delivered mathematical text is the author's own, copied from the pinned Stacks Project chapter file pinned-original-source.tex; no statement or proof step was rewritten, reordered or removed.\par\medskip
\noindent\textit{Editorial note.} Following the multi-claim convention for this source, source\_url is the source-level URL https://stacks.math.columbia.edu/ and the source label lemma-extend-across is carried in claim\_key and here rather than in a per-claim tag URL.\par\medskip
\noindent\textit{Editorial note.} The preamble is the chapter's shared Stacks Project preamble, retained as sources/stacks-morphisms/retained-preamble.tex: a line-for-line excerpt of the shared Stacks preamble staged with this snapshot. No macro was invented or redefined, and this card adds no preamble text of its own.\par\medskip
\noindent\textit{Editorial note.} Statement span 10424-10438 and proof span 10440-10473 come from the halfA pool card operations/pool-cards/stacks-morphisms-halfA.json and were re-read in the pinned chapter file before this card was written.\par\medskip
\noindent\textit{Editorial note.} The proof contains an xymatrix square displaying the morphism from the spectrum of the local ring extending the one from the spectrum of its fraction field; it is copied verbatim and the retained chapter preamble loads Xy-pic.\par\medskip
\noindent\textit{Editorial note.} The proof contains the author's parenthesis '(small detail omitted)' for the harmless replacement of the source by an open neighbourhood, and a reference to the author's earlier lemma producing a morphism from a prescribed germ; both are preserved.\par\medskip
\noindent\textit{Editorial note.} 4 cross-references occur inside the counted spans; they are listed in external\_references as navigation pointers only.\par\medskip
\noindent\textit{Allowed standard prerequisites (editorial).}\begin{itemize}
\item The same-chapter valuative criterion of properness, in the form producing a morphism from the spectrum of a valuation ring compatible with a given morphism from its fraction field. Cited at proof lines 10443-10444. Produces the morphism on the spectrum of the local ring at the point, which is the candidate germ of the extension.
\item The lemma producing an open neighbourhood and a morphism with prescribed germ, cited from the same chapter. Cited at proof lines 10449-10452. Spreads the germ at the point out to an actual morphism on a neighbourhood.
\item The equalizer of two morphisms to a separated scheme over the base is a closed subscheme, cited from the schemes chapter. Cited at proof lines 10455-10457. Shows that the two morphisms agree on a closed subscheme of the overlap, which contains the generic point.
\item For every prime of a scheme there is an affine open containing it with the property that the ring of global functions injects into the local ring at that prime, cited from the properties chapter. Cited at proof lines 10427-10430. Lets the proof assume the source affine and integral with the local ring at the point a valuation ring whose fraction field is the residue field of the generic point.
\end{itemize}\medskip
\ifdefined\hypertarget\hypertarget{sourcepointer-1}{}\fi
\noindent\textit{Source reference.} Original lemma 0BX5: focused-seeds/AGP/sources/stacks-snapshot-20261011/tags.txt, source lines 14968--14968. Same-chapter reference used in the retained author text at source line(s) 10451. It is supplied only as a navigation pointer so that the retained text compiles; the counted item delivers the author's own statement and the author's complete proof, and no part of the referenced item is copied in or substituted for the counted argument.\par\medskip
\ifdefined\hypertarget\hypertarget{sourcepointer-2}{}\fi
\noindent\textit{Source reference.} Original lemma 0BX6: focused-seeds/AGP/sources/stacks-snapshot-20261011/tags.txt, source lines 14969--14969. Same-chapter reference used in the retained author text at source line(s) 10461, 12656. It is supplied only as a navigation pointer so that the retained text compiles; the counted item delivers the author's own statement and the author's complete proof, and no part of the referenced item is copied in or substituted for the counted argument.\par\medskip
\ifdefined\hypertarget\hypertarget{sourcepointer-3}{}\fi
\noindent\textit{Source reference.} Original lemma 0BX3: focused-seeds/AGP/sources/stacks-snapshot-20261011/tags.txt, source lines 14966--14966. Same-chapter reference used in the retained author text at source line(s) 10443. It is supplied only as a navigation pointer so that the retained text compiles; the counted item delivers the author's own statement and the author's complete proof, and no part of the referenced item is copied in or substituted for the counted argument.\par\medskip
\ifdefined\hypertarget\hypertarget{sourcepointer-4}{}\fi
\noindent\textit{Source reference.} Original lemma 01KM: focused-seeds/AGP/sources/stacks-snapshot-20261011/tags.txt, source lines 2055--2055. Same-chapter reference used in the retained author text at source line(s) 1222, 10467, 12788, 12834. It is supplied only as a navigation pointer so that the retained text compiles; the counted item delivers the author's own statement and the author's complete proof, and no part of the referenced item is copied in or substituted for the counted argument.\par\medskip
\begin{lemma}
\label{lemma-extend-across}
Let $S$ be a scheme. Let $X$, $Y$ be schemes over $S$. Let $x \in X$.
Let $U \subset X$ be an open and let $f : U \to Y$ be a morphism over $S$.
Assume
\begin{enumerate}
\item $x$ is in the closure of $U$,
\item $X$ is reduced with finitely many irreducible components or
$X$ is Noetherian,
\item $\mathcal{O}_{X, x}$ is a valuation ring,
\item $Y \to S$ is proper
\end{enumerate}
Then there exists an open $U \subset U' \subset X$ containing
$x$ and an $S$-morphism $f' : U' \to Y$ extending $f$.
\end{lemma}

\begin{proof}
It is harmless to replace $X$ by an open neighbourhood of $x$ in $X$
(small detail omitted). By Properties, Lemma
\ref{properties-lemma-ring-affine-open-injective-local-ring}
we may assume $X$ is affine with
$\Gamma(X, \mathcal{O}_X) \subset \mathcal{O}_{X, x}$.
In particular $X$ is integral with a unique generic point $\xi$
whose residue field is the fraction field $K$ of the
valuation ring $\mathcal{O}_{X, x}$.
Since $x$ is in the closure of $U$ we see that $U$ is not
empty, hence $U$ contains $\xi$. Thus by the valuative criterion
of properness (Lemma \ref{lemma-characterize-proper})
there is a morphism $t : \Spec(\mathcal{O}_{X, x}) \to Y$
fitting into a commutative diagram
$$
\xymatrix{
\Spec(K) \ar[d]_\xi \ar[r] & \Spec(\mathcal{O}_{X, x}) \ar[d]_t \\
U \ar[r]^f & Y
}
$$
of morphisms of schemes over $S$. Applying
Lemma \ref{lemma-morphism-defined-local-ring}
with $y = t(x)$ and $\varphi = t^\sharp_x$ we obtain an open
neighbourhood $V \subset X$ of $x$ and a morphism $g : V \to Y$
over $S$ which sends $x$ to $y$ and such that $g^\sharp_x = t^\sharp_x$.
As $Y \to S$ is separated, the equalizer $E$ of $f|_{U \cap V}$
and $g|_{U \cap V}$ is a closed subscheme of $U \cap V$, see
Schemes, Lemma \ref{schemes-lemma-where-are-they-equal}.
Since $f$ and $g$ determine the same morphism $\Spec(K) \to Y$
by construction we see that $E$ contains the generic point
of the integral scheme $U \cap V$. Hence $E = U \cap V$ and
we conclude that $f$ and $g$ glue to a morphism $U' = U \cup V \to Y$
as desired.
\end{proof}

\end{document}
